\subsection{数值域}

定义 2.　设 E 为复 Hilbert 空间，u 为 E 的自同态。我们把 \(u\) 的数值域称为形如 \(\langle x|u(x)\rangle\) 的复数所成的集合，其中 x 取遍 E 的单位球面；把它记作 \(\iota(u)\) 。

\(u^{*}\) 的数值域是 \(\iota(u)\) 在复共轭下的像。对所有复数 \(\lambda\) 与 \(\mu\) ，\(\lambda u + \mu 1_{\mathrm{E}}\) 的数值域等于 \(\lambda \iota(u) + \mu\) 。

命题 6.　设 E 为复 Hilbert 空间，u 为 E 的自同态。

\begin{enumerate}
\def\labelenumi{\alph{enumi})}
\item
  u 的特征值所成的集合包含于 \(\iota(u)\) ；
\item
  u 的谱包含于 \(\iota(u)\) 在 C 中的闭包。
\end{enumerate}

设 \(\lambda\) 是 \(u\) 的特征值，\(x \in \mathrm{E}\) 为满足 \(u(x) = \lambda x\) 的非零向量。把 \(x\) 换成 \(x / \| x\|\) ，我们就可以假设 \(\| x \| = 1\) 。这时 \(\langle x | u(x) \rangle = \lambda\) ，因而 \(\lambda \in \iota(u)\) 。

我们证明 \(b\) ）。考虑 \(u - \lambda 1_{\mathrm{E}}\) ，便可归结为证明：若 0 属于 \(u\) 的谱，则 0 是 \(\iota(u)\) 的闭包点。

首先假设存在实数 \(c > 0\) ，使得 \(\| u(x)\| \geqslant c\) 对 E 中所有范数为 1 的 \(x\) 成立。这时自同态 \(u\) 是单射且其像是闭的（I，第 107 页，引理 8）。由于按假设它不可逆，故它不是满射。于是 \(u^{*}\) 的核的正交补不等于 E（EVT，V，第 41 页，命题 4），这就证明 \(u^{*}\) 的核非零。因此 0 属于 \(\iota(u^{*})\) ，从而属于 \(\iota(u)\) 。

若前面的假设不成立，则对每个整数 \(n \geqslant 1\) ，存在 E 中范数为 1 的向量 \(x_{n}\) ，满足 \(\|u(x_{n})\| \leqslant 1/n\) 。这时我们有 \(|\langle x_{n}|u(x_{n})\rangle| \leqslant 1/n\) ，由此推出 0 属于 \(\iota(u)\) 的闭包。

命题 7（Hausdorff–Toeplitz 定理）

设 E 为复 Hilbert 空间，\(u \in \mathcal{L}(\mathrm{E})\) 。数值域 \(\iota(u)\) 是 C 的一个凸子集。

为证明这个命题，我们需要两个引理。

引理 5.　设 E 为 2 维复 Hilbert 空间。给 E 的 Hermite 自同态所成的实向量空间 \(\mathcal{L}(\mathrm{E})_h\) 赋以准 Hilbert 范数 \(u \mapsto \mathrm{Tr}(u^* u)^{1/2}\) 。E 的秩 1 正交投影所成的集合是以 \(1/\sqrt{2}\) 为半径、以 \(\frac{1}{2} 1_{\mathrm{E}}\) 为心的球面 S，它位于 \(\mathcal{L}(\mathrm{E})_h\) 中由迹为 1 的自同态组成的 3 维仿射子空间内。

设 F 为 \(\mathcal{L}(\mathrm{E})_h\) 中由迹为 1 的元素组成的实仿射子空间。E 的秩 1 正交投影都属于 F（引理 3，(ii)）。

设 \(u \in \mathrm{F}\) 。我们有 \(\| u - \frac{1}{2} 1_{\mathrm{E}} \|^{2} = \mathrm{Tr}(u^{2} - u + \frac{1}{4}) = \mathrm{Tr}(u^{2}) - \frac{1}{2}\) 。因此，\(u \in \mathrm{S}\) 的充分必要条件是 \(\mathrm{Tr}(u^{2}) = 1\) 。由于 \(2\det(u) = \mathrm{Tr}(u)^{2} - \mathrm{Tr}(u^{2}) = 1 - \mathrm{Tr}(u^{2})\) ，这个条件等价于 \(\det(u) = 0\) 。于是由 Hamilton–Cayley 定理（A，III，第 107 页，命题 20），我们有 \(u \in \mathrm{S}\) 的充分必要条件是 \(u^{2} - u = 0\) ，这意味着 \(u\) 是秩 1 的 Hermite 投影（同前），由此即得结论。

引理 6.　设 E 为实赋范向量空间，F 为实向量空间，u: E → F 为一个非单射的仿射映射。设 B 是 E 中的一个球，S 是相应的球面。我们有 \(u(S) = u(B)\) ，特别地，\(u(S)\) 是凸的。

我们可以归结到 \(u\) 是线性映射且 B 是 E 的单位球的情形。我们有 \(u(S) \subset u(B)\) 。反之，设 \(x \in B\) ，并设 \(y\) 是 \(\text{Ker}(u)\) 中的非零元素。连续映射 \(t \mapsto \|x + ty\|\) （它把 \(\mathbf{R}\) 映到 \(\mathbf{R}_+\) 中）的像是一个包含实数 \(\|x\| \leqslant 1\) 的无界区间。因此存在 \(t \in \mathbf{R}\) ，使得 \(\|x + ty\| = 1\) 。这时我们有 \(x + ty \in S\) 与 \(u(x + ty) = u(x)\) ，因而 \(u(x) \in u(S)\) 。

我们来证明命题 7。设 x 与 y 是 E 的单位球面的元素；我们证明以 \(\langle x|u(x)\rangle\) 与 \(\langle y|u(y)\rangle\) 为端点的线段包含在 u 的数值域中。

设 F 为 E 的由 x 与 y 生成的子空间。若 \(\dim(\mathrm{F}) = 1\) ，则有 \(\langle x|u(x)\rangle = \langle y|u(y)\rangle\) ，断言得证。否则，我们有 \(\dim(\mathrm{F}) = 2\) ；设 p 是 E 的以 F 为像的正交投影，并用 \(u_{F}\) 表示由 \(x \mapsto p(u(x))\) 所给出的 F 的自同态。由于 p 是 Hermite 的（I，第 133 页，引理 3），我们有 \(\langle z|u_{\mathrm{F}}(z)\rangle = \langle z|u(z)\rangle\) 对所有 \(z \in F\) 成立，从而 \(\iota(u_{\mathrm{F}}) \subset \iota(u)\) 。因此我们可以假设 E = F。

对 E 的每个元素 \(z\) ，用 \(v_{z}\) 表示由 \(t \mapsto \langle z|t\rangle z\) 所定义的 E 的 Hermite 自同态；我们有 \(\langle z|u(z)\rangle = \mathrm{Tr}(u \circ v_{z})\) 。当 \(z\) 取遍 E 的单位球面时，\(v_{z}\) 描出 E 的秩 1 正交投影所成的集合，即 \(\mathcal{L}(\mathrm{E})\) 的由迹为 1 的 Hermite 自同态组成的实仿射子空间 V 中的一个球面 S（引理 5）。因此，u 的数值域就是诸 \(\mathrm{Tr}(u \circ v)\) （\(v \in \mathrm{S}\)）所成的集合。映射 \(v \mapsto \mathrm{Tr}(u \circ v)\) 把 V 映到 \(\mathbf{C}\) 中，它是线性的。由于 \(\dim_{\mathbf{R}}(\mathrm{V}) = 3 > \dim_{\mathbf{R}}(\mathbf{C})\) ，它不是单射；于是由引理 6 便知 \(\iota(u)\) 是凸的。

